\section{Introduction}\label{sec:introduction}

Open-source scientific software has lowered the cost of applying sophisticated
statistical methods. Libraries such as scikit-learn
\citep{pedregosa-etal-2011}, statsmodels \citep{seabold-perktold-2010}, dcor
\citep{ramos-carreno-torrecilla-2023} and lifelines \citep{davidson-pilon-2019}
make machine-learning, econometric, distributional and survival methods
available through consistent interfaces, and each was introduced with a short
paper describing its design.

Implementing a statistical method and justifying it are nevertheless different
activities. A procedure combines a target quantity, a model and its
assumptions, a numerical algorithm, an uncertainty estimate and a reference
distribution. A converged optimizer establishes none of the inferential claims
made from its output, and an algebraic identity between two formulas does not
make their floating-point evaluations interchangeable. These distinctions
matter more as software becomes easier to extend. Automatic differentiation and
compilation through JAX \citep{bradbury-etal-2018} make sophisticated numerical
programs practical, and AI-assisted programming reduces the effort of producing
candidate implementations. Neither change makes it easier to establish what an
implementation guarantees.

Proof assistants offer a complementary capability. Lean~4
\citep{de-moura-ullrich-2021} checks mathematical developments mechanically,
Mathlib \citep{mathlib-2020} supplies reusable foundations, and verified
compiler work shows how proofs can establish that well-specified program
rewrites preserve meaning \citep{bhat-etal-2024}. A proof, however, establishes
a stated proposition under its hypotheses. It does not by itself certify a
Python function, a compiled kernel or the finite-sample behaviour of a
statistical procedure.

STATADIC is a statistical computing library built around these boundaries. Its
contribution is not a claim to replace established packages, nor a collection
of universally validated estimators. It is an architecture in which statistical
components are reusable while their mathematical meaning, numerical behaviour
and supporting evidence stay distinguishable. The paper makes three
contributions:
\begin{enumerate}
  \item a separation of statistical responsibilities, from data representation
    to reporting, that lets methods share components without erasing the
    assumptions that govern them (\cref{sec:design,sec:workflow});
  \item a layered treatment of assurance, in which machine-checked
    mathematics, implementation correspondence, numerical conformance,
    statistical qualification and performance are recorded as separate claims
    (\cref{sec:assurance});
  \item evidence produced from recorded source by a single script, including
    results that are unfavourable to procedures STATADIC implements
    (\cref{sec:evidence}).
\end{enumerate}
\Cref{sec:discussion} discusses costs and limitations.
